\section{Concluding Remarks and Future Plans}\label{sec:conclusions}

We reported on an experimental study exploring the suitability of LLMs as substitutes for human participants in evaluating the quality of language vocabularies according to a proposed methodology, Peira. We constructed a data collection instrument intended for humans, reproduced it in the form of an LLM prompt containing the same information, and administered each to a group of humans and a group of LLMs, respectively. We found that LLM performance level is too distant from that of humans to allow the former to be used as a substitute for the latter, and that performance tends to improve as model size increases in number of parameters. 

The result opens many research questions and opportunities for future investigation. One such question pertains to finding explanations for the result, i.e., why humans perform worse than LLMs even when both groups are given the exact same information and both arguably have the preparation to comprehend that material. One explanation points to the quality of the voice and visuals in the training video or the layout of the instrument, which participants may find confusing or unstimulating. An additional, perhaps more credible, concern is that the online participants, despite having self-reported information technology education at the university or college level, may lack the specific interest or background in matters pertinent to iStar-DT. However, one can counter-argue that concepts such as \textit{goal}, \textit{effect}, \textit{wants-to} are readily comprehensible by an even wider population than the one from which our participants were sampled. Regardless, such concerns may be construed, respectively, as \textit{internal} and \textit{external validity} threats to our study and will need follow-up investigations to be resolved. 

Furthermore, the observed human vs. LLM performance gap inspires us to ask what level, if any, of training effort, intensity, detail, and care suffices for motivated humans to reach LLM-level performance. This, in turn, leads us to rethink how processes like Peira are to be utilized. Assuming that the performance of powerful LLMs offers a vocabulary comprehension standard, then Peira, when applied to humans, can be used to measure the \textit{learnability} of the vocabulary or the quality of the training process, rather than the quality of the vocabulary itself. In the iStar-DT case, for example, we could assume that its vocabulary is as well-designed as LLM ratings imply it is, and attribute the discrepancy with human data to low learnability. This would imply that the designers may either call the vocabulary successful but caution that it needs extensive training (surely more than what LLMs require), or re-engineer it to make it more learnable, which may, however, imply expressiveness-suppressing decisions, such as reducing the size of the vocabulary through concept abstraction. Nevertheless, the premise of the above reasoning, i.e., that LLMs offer a standard of comprehension of a language that any future human user can attain if they just receive enough training, needs itself to be tested and qualified. We plan to design and conduct studies that clarify if such reasoning is valid. 

We are also interested in acquiring a more complete picture of how LLMs can assist with the assessment of vocabulary quality. For example, self-reported metrics such as of semantic overlaps, and of construct deficits and excesses, which Peira elicits from humans, can be investigated when provided by LLMs vis-\`{a}-vis their relevance and validity, through, for example, expert assessment. We also plan to explore whether qualitative feedback and advice from LLMs can be useful in improving the language. 

As a final remark, in conceptualizing future work on modeling language quality assessment in general, it is relevant to acknowledge that in a world where LLM agents are also, alongside humans, producers and users of models in practical production environments, LLMs become consumers of quality as well. They are, hence, additional stakeholders to satisfy in designing the language and its training process and material. Developing systematic empirical evaluation practices tailored to LLMs is, thus, another direction for future investigation.



